

\paragraph{Begin another period with recorded balances and unresolved findings.}
A conditional simulation can need another period before all the preceding
period's accounting questions have been resolved. For example, the scenario
above has cash, an unpaid current-tax liability, lease rights and a recorded
profit, but its deferred-tax assessment and the completeness of its historical
inputs remain open. Those recorded quantities are sufficient to construct a
provisional opening for selected subsequent transactions. They are insufficient
to declare that the preceding annual financial statements are complete.

The explicit model \texttt{provisional\_recorded\_balance\_carry\_v1} makes
that distinction part of execution. The request identifies the complete prior
reassessment, a unique carry event and the next period's start and exclusive
end. The prior current-tax assessment and loss history must reach the preceding
financial endpoint, and the next period must start at that endpoint. The prior
closing date must be its last included day. No calendar year or assessment of
missing transactions is inferred from those dates. Separate findings describe
event coverage, deferred-tax assessment, source admission and accounting
admission. Unknown and negative findings remain visible in the resulting state;
affirmative caller assumptions also leave annual readiness unestablished.

The transfer follows the book's debit-positive convention. Let $P$ denote the
recorded profit, $R$ the signed retained-earnings balance, and $b_k$ each
signed account balance immediately before the boundary. Let $\mathcal I$
contain the temporary income and expense accounts. The new opening satisfies
\[
 P=-\sum_{k\in\mathcal I}b_k,\qquad
 R^{\mathrm{open}}=R-P,\qquad
 b_k^{\mathrm{open}}=0\quad(k\in\mathcal I).
\]
Every other account keeps its recorded amount. A profit therefore increases the
credit balance in retained earnings; a loss reduces it. The transfer journal
sums to zero without a balancing plug. In the mixed scenario, recorded profit
is $-1/4$ and retained earnings is $-12$, so the opening signed balance is
$-12-(-1/4)=-47/4$. Cash stays $721/2$ and the unpaid tax liability stays
$19/4$. Positive-profit and zero-profit cases separately check the same sign
convention.

The boundary changes period accumulators without discarding economic memory.
Return components, recovered lots including disposed identities, lease rights,
measurement dates, loss quantities and lifetime event identifiers survive
exactly. The prior result retains its cash flows and tax adjustments together
with hashes of the closing book and report. The next period begins with empty
cash-flow and tax-adjustment lists, so an old receipt or deduction cannot be
counted again. Opening deferred-tax allocations copy the recorded asset and
liability; that copying does not assess their recognition or measurement. The
new period's assessment cursor denotes an empty interval and its current tax is
explicitly unassessed.

Subsequent statements use the actual carried opening. Paying the opening tax
liability of $19/4$ reduces cash from $721/2$ to $1423/4$, eliminates that
liability and leaves current-period tax expense at zero. In the indirect cash
calculation, the reduction in the opening payable explains the operating cash
outflow. It must not be charged to profit again. If depreciation of 10 follows,
recorded profit is $-10$, assets are $2743/4$, liabilities are 104 and equity is
$2327/4$. Depreciation is added back in the indirect bridge, leaving operating
cash flow at $-19/4$. The balance sheet, equity movement and direct and indirect
cash calculations reconcile to the same opening and closing state.

A second example carries an unsettled return component instead. Settlement
pays 48, releases a refund liability of 40 and a recovery asset of 16, and
recognizes recovered inventory of 20. Its cost-of-sales credit of 4 and selling
expense of 8 reduce profit by 4. Selling that inventory for 30 then earns 10,
so the two transactions together earn 6 and reduce cash by 18. Starting from
cash of 500, the following-period cash is 482. The indirect calculation gives
$6+16-40=-18$. The new return tax adjustment of $-24$ is retained for a later
assessment; it is not absorbed into an asserted tax-complete closing balance.

Ordered execution retains the completed prefix if a later event fails. A
payment still needs both an unpaid liability and actual shared cash; a reused
identity or an event outside the next period is refused. Applying a result
recomputes the entire predecessor and subsequent prefix and compares the
complete result, including unresolved findings and archived history. The
checked implementation covers this provisional boundary and selected return,
lease and tax-payment operations. New-period tax assessment, repeated carry,
complete event coverage and deferred-tax recognition require their own methods
and evidence before the scenario can support a complete annual close.
